\section{Lean4 Formalization}
\label{app:lean}

This appendix describes the formal verification artifact at a level intended for
paper readers. It is not a theorem-name index. The theorem-by-theorem mapping,
assumption crosswalk, and module-level navigation notes live in the repository
docs under \texttt{lean3/docs/}, where they can evolve without making the paper
appendix unreadable.

\subsection{Artifact Summary}

The Lean sources live under \texttt{lean3/}. The active development currently
contains more than 50{,}000 lines across 123 Lean files. The code is organized
into five main parts:
\begin{itemize}[nosep]
    \item \textbf{OPT}: the core Oracle-Preserving Summarization theory, including local laws, preference-learning equivalence, gap bounds, theorem-backed reduction, measurement-error extensions, and adaptive-tree optimizer transfer.
    \item \textbf{DSL}: the design-based statistical layer, including IPW theory, audit certificates, nonclassical expectation mismatch, and TreePO end-to-end guarantees.
    \item \textbf{CLT}: probability infrastructure used for asymptotic arguments.
    \item \textbf{Econometrics}: supporting results for identification, regression, panel methods, and related foundations.
    \item \textbf{ML}: general supervised-learning primitives used by the broader theory stack.
\end{itemize}

The formalization builds on Mathlib and a companion probability library. The
active build contains one axiom and no \texttt{sorry} placeholders in the
imported theorem stack.

\subsection{What Is Formally Verified}

At a high level, the formalization verifies the full chain used in the paper:
\begin{enumerate}[nosep]
    \item local laws imply one-pass, schedule-invariant, and multi-round preservation;
    \item exact preservation implies DPO, GRPO-PL, and GRPO-RL equivalence;
    \item approximate preservation implies quantitative gap bounds for those same objectives;
    \item theorem-backed reduction can be exact or approximate, can be direct or codec-mediated, and can be analyzed under measurement error;
    \item stochastic adaptive tree objectives inherit optimizer-transfer guarantees, including cases where the oracle is imperfect and cases where certification holds only on a high-probability event;
    \item clustered confidence intervals, calibration radii, and oracle-measurement envelopes feed directly into the final \texttt{computeDSLBound} object through machine-checked event-to-bound theorems;
    \item the audit and IPW layers propagate these guarantees through the TreePO / DSL certificate stack.
\end{enumerate}

The same development also formalizes the bridge to classical mergeable
summaries, proves that merge idempotence alone is not enough for OPS
idempotence, and isolates the exact additional assumptions needed for
theorem-backed learned reductions.

\subsection{Single Axiom}

The active theorem stack uses one formal axiom: expected group loss is
Lipschitz in oracle distance. This is the precise point where the paper appeals
to standard random-utility reasoning. The motivation is that preference losses
based on rankings are often not pointwise Lipschitz because rankings can change
discontinuously, but they \emph{are} expected-Lipschitz once one integrates over
a continuous noise model such as Gumbel or Gaussian perturbations. That is the
economically standard move behind Plackett--Luce and related discrete-choice
models, and it is the only non-derived ingredient in the active formal proof
stack.

\subsection{Human-Readable Proofs vs. Repository Maps}

The paper appendix follows a simple rule:
\begin{itemize}[nosep]
    \item if a mathematical claim appears in the appendix, it should have a full human-readable proof there;
    \item low-level Lean identifiers, paper-to-Lean mappings, and assumption crosswalks belong in repository documentation rather than in the paper itself.
\end{itemize}

Accordingly, readers who want the exact code-level correspondence should use the
standalone repository files:
\begin{itemize}[nosep]
    \item \texttt{lean3/docs/CORE\_PROOFS.md}: paper proofs reconstructed step by step alongside the Lean structure;
    \item \texttt{lean3/docs/PAPER\_TO\_LEAN\_MAP.md}: theorem-by-theorem mapping from paper labels to Lean names and files;
    \item \texttt{lean3/docs/ASSUMPTION\_CROSSWALK.md}: paper assumptions and their Lean counterparts;
    \item \texttt{lean3/docs/ADAPTIVE\_TREE\_OPTIMIZER\_TRANSFER.md}: the expected-tree and high-probability optimizer-transfer results.
\end{itemize}

For code readers, \texttt{lean3/FormalProofs/OPT/MainTheorems.lean} is the
curated Lean entry point.

\subsection{Coverage Boundary}

The formalization covers the mathematical guarantees of the system, not every
engineering detail of the implementation. In particular:
\begin{itemize}[nosep]
    \item \textbf{Formalized}: local laws, reduction guarantees, preference-learning objectives, measurement-error variants, adaptive-tree optimization theorems, audit/sample-complexity results, and the IPW / DSL certificate layer.
    \item \textbf{Not formalized}: low-level LLM I/O, batching, HTTP behavior, serialization, GPU orchestration, and the operational details of the Python runtime.
\end{itemize}

This boundary is deliberate. The proofs certify the mathematical contracts that
the software is supposed to satisfy. They do not attempt to prove the correctness
of every line of orchestration code.

\subsection{Build Instructions}

To verify the full formalization:

\begin{lstlisting}
cd lean3
lake build FormalProofs
\end{lstlisting}

For the core OPS entry point only:

\begin{lstlisting}
cd lean3
lake build FormalProofs.OPT.MainTheorems
\end{lstlisting}

The repository contains CI configuration for rebuilding these targets
automatically on revision changes.
